\documentclass[12pt,a4paper]{article}

\usepackage[utf8]{inputenc}
\usepackage[margin=1in]{geometry}
\usepackage{titlesec}
\usepackage{booktabs}
\usepackage{tabularx}
\usepackage{xcolor}
\usepackage{listings}
\usepackage{enumitem}
\usepackage{hyperref}

\hypersetup{
    colorlinks=true,
    linkcolor=black,
    urlcolor=blue,
    citecolor=black
}

\definecolor{codebg}{rgb}{0.96, 0.96, 0.96}
\definecolor{codeframe}{rgb}{0.85, 0.85, 0.85}
\lstset{
    backgroundcolor=\color{codebg},
    frame=single,
    rulecolor=\color{codeframe},
    basicstyle=\ttfamily\footnotesize,
    breaklines=true,
    numbers=left,
    numberstyle=\tiny\color{gray},
    tabsize=2,
    showstringspaces=false
}

\begin{document}

\begin{titlepage}
    \centering
    \vspace*{1.5cm}
    {\Large \textbf{A PROJECT REPORT ON}}\\[0.8cm]
    {\Huge \textbf{GymLog: Fitness Tracking Mobile Application}}\\[1.2cm]
    {\large Submitted in partial fulfillment of the requirements\\for the degree of}\\[0.4cm]
    {\large \textbf{Bachelor of Technology / Computer Science}}\\[2.5cm]
    \begin{minipage}{0.45\textwidth}
        \begin{flushleft} \large
            \emph{Submitted By:}\\
            \textbf{Siddhesh Pandey}\\
            Department of Computer Science
        \end{flushleft}
    \end{minipage}
    \begin{minipage}{0.45\textwidth}
        \begin{flushright} \large
            \emph{Project Guide:}\\
            \textbf{[Faculty / Guide Name]}\\
            Assistant Professor
        \end{flushright}
    \end{minipage}
    \vfill
    {\large \textbf{[Your College / University Name]}}\\[0.2cm]
    {\large Academic Year 2026--2027}
\end{titlepage}

\tableofcontents
\newpage

\section{Executive Summary}
\textbf{GymLog} is an intuitive, distraction-free mobile fitness tracking application inspired by modern workout loggers such as Hevy. Built using the \textbf{Flutter} framework for the mobile interface and \textbf{Node.js}, \textbf{Express}, and \textbf{MongoDB Atlas} for backend cloud services, GymLog provides an eye-friendly dark theme optimized specifically for strength and resistance training.

\section{Project Scope and Objectives}
\subsection{Project Objectives}
\begin{enumerate}[noitemsep]
    \item To build an intuitive, dark-themed mobile user interface that minimizes cognitive friction during high-intensity gym sessions.
    \item To implement a secure RESTful backend API using JSON Web Tokens (JWT) for stateless authentication and bcrypt for password encryption.
    \item To provide interactive live workout session tracking with automated countdown rest timers between sets.
    \item To unify performance telemetry across multiple routines so that an exercise retains its historical progress regardless of where it is performed.
    \item To provide social accountability by allowing athletes to publish workout summaries with photos and interact through likes.
\end{enumerate}

\subsection{In-Scope Features}
\begin{itemize}[noitemsep]
    \item \textbf{User Authentication}: Registration, login, and secure local token storage.
    \item \textbf{Exercise Database}: 80+ default movements categorized by muscle group and equipment.
    \item \textbf{Routine Management}: Complete CRUD for custom workout routines.
    \item \textbf{Active Logging}: Set check-off, weight/repetition inputs, dynamic rest timer controls (+30s, Skip).
    \item \textbf{Progress Analytics}: Time-series charts for workout volume, max weights, and body weight logs.
    \item \textbf{Workout History}: Detailed inspection, modification, and deletion of past workout sessions.
    \item \textbf{Community Feed}: Public timeline of completed workouts with captions, media photos, and likes.
\end{itemize}

\section{System Architecture}
GymLog employs a Three-Tier Client-Server Architecture:
\begin{enumerate}
    \item \textbf{Presentation Tier (Mobile Client)}: Implemented in Flutter with Material 3 Dark theme, utilizing Provider for state management.
    \item \textbf{Application Tier (API Gateway)}: Implemented using Node.js and Express with JWT authentication middleware.
    \item \textbf{Data Tier (Database)}: A MongoDB Atlas cloud database storing JSON-like documents via Mongoose schemas.
\end{enumerate}

\section{Technology Stack}
\begin{table}[h!]
    \centering
    \caption{Summary of Technologies Used}
    \begin{tabularx}{\textwidth}{l l X}
        \toprule
        \textbf{Component} & \textbf{Technology} & \textbf{Role in System} \\
        \midrule
        \textbf{Frontend} & Flutter (Dart 3.x) & Cross-platform framework for Android and iOS \\
        & Provider & Reactive state management \\
        & fl\_chart & Interactive charting library for progress tracking \\
        & flutter\_secure\_storage & Hardware-backed keychain/keystore token storage \\
        & flutter\_dotenv & Environment variable configuration loader \\
        \midrule
        \textbf{Backend} & Node.js 20+ & Asynchronous server runtime \\
        & Express.js 5.x & RESTful API routing and middleware framework \\
        & Mongoose 9.x & Object Data Modeling (ODM) for MongoDB \\
        & jsonwebtoken & Stateless user authentication mechanism \\
        & bcryptjs & One-way password hashing algorithm \\
        & cors & Cross-Origin Resource Sharing handling \\
        \midrule
        \textbf{Database} & MongoDB Atlas & Managed cloud NoSQL document database \\
        \bottomrule
    \end{tabularx}
\end{table}

\section{Installation and Setup Guide}
\subsection{Prerequisites}
Node.js (v18+), Flutter SDK (v3.19+), and access to MongoDB Atlas.

\subsection{Backend Setup}
\begin{lstlisting}[language=bash]
cd gymlog_backend
npm install
npm run seed  # Populates demo athlete profile (siddhesh@gymlog.app / password123)
npm start
\end{lstlisting}

\subsection{Frontend Setup}
\begin{lstlisting}[language=bash]
cd gymlog
flutter pub get
flutter run
\end{lstlisting}

\section{API Reference}
\begin{itemize}[noitemsep]
    \item \textbf{POST /api/auth/register}: Creates a new user profile.
    \item \textbf{POST /api/auth/login}: Verifies credentials and returns a JWT.
    \item \textbf{GET /api/routines} \& \textbf{POST /api/routines}: Read and create workout routines.
    \item \textbf{POST /api/workouts}: Logs a completed workout session with sets, reps, and media.
    \item \textbf{GET /api/workouts/exercise-history/:name}: Cross-routine progression history.
    \item \textbf{GET /api/bodyweight} \& \textbf{POST /api/bodyweight}: Body weight tracking.
    \item \textbf{GET /api/feed} \& \textbf{POST /api/feed/:id/like}: Community feed interactions.
\end{itemize}

\section{Conclusion}
The GymLog project successfully demonstrates the design and implementation of a full-stack mobile application for fitness tracking, combining a high-performance Flutter mobile client with a scalable Node.js and MongoDB cloud backend.

\end{document}
